\begin{proposition}[6.8.3]
\label{IV.2.s6.8.3-fr}
Para un morfismo $f$ de preesquemas localmente noetherianos, sea $\boldsymbol{M}(f)$
una de las siguientes propiedades: ser de Cohen-Macaulay, $(S_n)$, regular,
$(R_n)$, normal o reducido.
\begin{enumerate}
\item[\textup{(i)}] Sean $X$, $Y$, $Z$ tres preesquemas localmente noetherianos,
$f:X\to Y$, $g:Y\to Z$ dos morfismos. Si $\boldsymbol{M}(f)$ y $\boldsymbol{M}(g)$ se cumplen,
se cumple $\boldsymbol{M}(g\circ f)$.
\item[\textup{(ii)}] Recíprocamente, si $f$ es sobreyectivo y si $\boldsymbol{M}(f)$ y $\boldsymbol{M}(g\circ f)$
se cumplen, entonces se cumple $\boldsymbol{M}(g)$.
\item[\textup{(iii)}] Sean $X$, $Y$, $Y'$ tres preesquemas localmente noetherianos,
$f:X\to Y$, $h:Y'\to Y$ dos morfismos; se pone $X'=X\times_Y Y'$,
$f'=f_{(Y')}:X'\to Y'$. Se supone que $f$ o $h$ es localmente de tipo
finito. Entonces, si se cumple $\boldsymbol{M}(f)$, también se cumple $\boldsymbol{M}(f')$; el recíproco
es válido cuando $h$ es fielmente plano.
\end{enumerate}
Las conclusiones de \textup{(i)} y \textup{(iii)} siguen siendo válidas cuando $\boldsymbol{M}$ es la propiedad
de ser suave y (en \textup{(iii)}) $h$ es cuasicompacto.
\end{proposition}

\begin{proof}
\begin{enumerate}
\item[(i)] Ya se sabe que si $f$ y $g$ son planos, también lo es
$u=g\circ f$, y que si $f$ y $g\circ f$ son planos y $f$ es sobreyectivo, $g$ es
plano (2.1.6 y 2.2.13). Por otra parte, para todo $z\in Z$, el morfismo
$f_z=f\otimes1_{\mathbf{k}(z)}:u^{-1}(z)\to g^{-1}(z)$ es plano (resp.
fielmente plano si $f$ lo es) y para todo $y\in g^{-1}(z)$,
$f_z^{-1}(y)$ es isomorfo a $f^{-1}(y)$ (\textbf I, 3.6.4). Si $\boldsymbol{M}(f)$ y $\boldsymbol{M}(g)$
se cumplen, se cumple por tanto $\boldsymbol{M}(g\circ f)$ en los casos en que $\boldsymbol{M}$ es la propiedad
de ser de Cohen-Macaulay o $(S_n)$, en virtud de (6.6.1). Por otra parte, sea
$k'$ una extensión finita de $\boldsymbol{k}(z)$; pongamos
$Y'_z=g^{-1}(z)\otimes_{\mathbf{k}(z)}k'$, $X'_z=u^{-1}(z)\otimes_{\mathbf{k}(z)}k'$, y
$f'_z=f_z\otimes1_{k'}:X'_z\to Y'_z$; el morfismo $f'_z$
\oldpage[IV-2]{152}
es plano
(resp. fielmente plano) y para todo $y'\in Y'_z$, la fibra $f_z'^{-1}(y')$
es isomorfa a $f^{-1}(y)\otimes_{\mathbf{k}(y)}\boldsymbol{k}(y')$, donde $y$ designa la imagen
de $y'$ en $g^{-1}(z)$ (\textbf I, 3.6.4). Cuando $\boldsymbol{M}$ es la propiedad de ser
regular, $(R_n)$, normal o reducido, la hipótesis de que $\boldsymbol{M}(f)$ y $\boldsymbol{M}(g)$
se cumplen implica que $Y'_z$ y cada una de las fibras $f_z'^{-1}(y')$ poseen,
para $y'\in Y'_z$, la propiedad correspondiente entre las propiedades \emph{c)}, \emph{d)},
\emph{e)}, \emph{f)} de (6.6.1); de (6.6.1, (i)) se deduce por tanto que $X'_z$ posee la misma
propiedad, de modo que se cumple $\boldsymbol{M}(g\circ f)$.
\item[(ii)] Recíprocamente, la hipótesis de que $\boldsymbol{M}(g\circ f)$ y $\boldsymbol{M}(f)$ se cumplen
y de que $f$ es sobreyectivo implica que $Y'_z$ tiene la propiedad correspondiente en
virtud de (6.6.1, (ii)), pues $f'_z$ es sobreyectivo para todo $z$; por tanto se cumple
$\boldsymbol{M}(g)$.
\item[(iii)] La primera afirmación se deduce inmediatamente de (6.8.2). Por otra parte,
si $h$ es fielmente plano y $f'$ es plano, $f$ es plano (2.4.1); como, con las
notaciones de (6.8.2), las propiedades $\boldsymbol{P}(f^{-1}(y),x)$ y
$\boldsymbol{P}(f'^{-1}(y'),x')$ son equivalentes (6.7.8), se ve que $\boldsymbol{M}(f)$ y $\boldsymbol{M}(f')$
son entonces equivalentes.
\end{enumerate}
Por último, la última afirmación de la proposición resulta de (1.4.3, (iii)) y
de (2.7.1, (iv)).
\end{proof}
